% GENERATED FILE. Edit canonical lessons, then run npm run generate:books.
% canonical-source-hash: fnv1a64:ce073b73a704ce3e
% canonical-lessons: GU-C414-purva, GU-C414-pascim, GU-C414-uttar, GU-C414-ultu, GU-C414-hokayantra

\chapter{East and West}
\label{ch:gu-east-and-west}
\hlchaptermodality{414}

\begin{chapteropening}
\textbf{By the end of this chapter:} \emph{I can say which way is north, east and west.}

Complete the last lesson of chapter 414: I can say which way is north, east and west.
\end{chapteropening}

\section[\texorpdfstring{\gu{પૂર્વ} (pūrva)}{પૂર્વ (pūrva)}]{\gu{પૂર્વ} (pūrva) — the east}

\label{lesson:GU-C414-purva}



\begin{quote}
Before the new one: say the Gujarati for to lead, then the Gujarati for to point out.
\end{quote}

\subsection*{You'll want to know: \gu{પૂર્વ}}
\textbf{\gu{પૂર્વ}} — \emph{pūrva} — \textquotedblleft{}the east\textquotedblright{}.

\begin{quote}
\gu{સૂરજ} \gu{પૂર્વમાં} \gu{ઊગે} \gu{છે}.
\end{quote}

\begin{quote}
\emph{sūraj pūrvamā̃ ūge che.}
\end{quote}

\begin{quote}
The sun rises in the east.
\end{quote}

A feminine noun.

\begin{cousinweb}
\textbf{\gu{પૂર્વ}} (\emph{pūrva}) is a tatsama, Sanskrit \emph{pūrva}, \textquotedblleft{}in front, before\textquotedblright{}, and so \textquotedblleft{}eastern\textquotedblright{} (Monier-Williams): facing the sunrise, the east is in front of you.
\end{cousinweb}

\begin{grammarlens}[title={What you've built}]
One more word for getting around.
\end{grammarlens}

\subsection*{Your turn}
\begin{itemize}
\raggedright
  \item \emph{Say it:} \emph{pūrva}
  \item \emph{Say it:} \emph{pūrva}, once more
  \item \emph{Say it:} \emph{sūraj pūrvamā̃ ūge che.} — The sun rises in the east.
  \item \emph{Recall:} say the Gujarati for to lead, then the Gujarati for to point out, then say \emph{pūrva} again
\end{itemize}

\begin{tcolorbox}[breakable,colback=teal!4,colframe=teal!35!black,title={Before you move on}]
What does \gu{પૂર્વ} mean? (\textquotedblleft{}The east\textquotedblright{}.) Say it once more.
\end{tcolorbox}
\section[\texorpdfstring{\gu{પશ્ચિમ} (paścim)}{પશ્ચિમ (paścim)}]{\gu{પશ્ચિમ} (paścim) — the west}

\label{lesson:GU-C414-pascim}



\begin{quote}
Before the new one: say the Gujarati for to point out, then the Gujarati for the east.
\end{quote}

\subsection*{You'll want to know: \gu{પશ્ચિમ}}
\textbf{\gu{પશ્ચિમ}} — \emph{paścim} — \textquotedblleft{}the west\textquotedblright{}.

\begin{quote}
\gu{સૂરજ} \gu{પશ્ચિમમાં} \gu{આથમે} \gu{છે}.
\end{quote}

\begin{quote}
\emph{sūraj paścimmā̃ āthme che.}
\end{quote}

\begin{quote}
The sun sets in the west.
\end{quote}

A feminine noun.

\begin{cousinweb}
\textbf{\gu{પશ્ચિમ}} (\emph{paścim}) is a tatsama, Sanskrit \emph{paścima}, \textquotedblleft{}behind, last\textquotedblright{}, and so \textquotedblleft{}western\textquotedblright{} (Monier-Williams): facing the sunrise, the west is behind you.
\end{cousinweb}

\begin{grammarlens}[title={What you've built}]
One more word for getting around.
\end{grammarlens}

\subsection*{Your turn}
\begin{itemize}
\raggedright
  \item \emph{Say it:} \emph{paścim}
  \item \emph{Say it:} \emph{paścim}, once more
  \item \emph{Say it:} \emph{sūraj paścimmā̃ āthme che.} — The sun sets in the west.
  \item \emph{Recall:} say the Gujarati for to point out, then the Gujarati for the east, then say \emph{paścim} again
\end{itemize}

\begin{tcolorbox}[breakable,colback=teal!4,colframe=teal!35!black,title={Before you move on}]
What does \gu{પશ્ચિમ} mean? (\textquotedblleft{}The west\textquotedblright{}.) Say it once more.
\end{tcolorbox}
\section[\texorpdfstring{\gu{ઉત્તર} (uttar)}{ઉત્તર (uttar)}]{\gu{ઉત્તર} (uttar) — the north}

\label{lesson:GU-C414-uttar}



\begin{quote}
Before the new one: say the Gujarati for the east, then the Gujarati for the west.
\end{quote}

\subsection*{You'll want to know: \gu{ઉત્તર}}
\textbf{\gu{ઉત્તર}} — \emph{uttar} — \textquotedblleft{}the north\textquotedblright{}.

\begin{quote}
\gu{અમારું} \gu{ગામ} \gu{શહેરની} \gu{ઉત્તરમાં} \gu{છે}.
\end{quote}

\begin{quote}
\emph{amārũ gām śahernī uttarmā̃ che.}
\end{quote}

\begin{quote}
Our village is to the north of the city.
\end{quote}

A feminine noun.

\begin{cousinweb}
\textbf{\gu{ઉત્તર}} (\emph{uttar}) is a tatsama, Sanskrit \emph{uttara}, \textquotedblleft{}upper, higher\textquotedblright{}, and so \textquotedblleft{}northern\textquotedblright{} (Monier-Williams). The same word also means an answer.
\end{cousinweb}

\begin{grammarlens}[title={What you've built}]
One more word for getting around.
\end{grammarlens}

\subsection*{Your turn}
\begin{itemize}
\raggedright
  \item \emph{Say it:} \emph{uttar}
  \item \emph{Say it:} \emph{uttar}, once more
  \item \emph{Say it:} \emph{amārũ gām śahernī uttarmā̃ che.} — Our village is to the north of the city.
  \item \emph{Recall:} say the Gujarati for the east, then the Gujarati for the west, then say \emph{uttar} again
\end{itemize}

\begin{tcolorbox}[breakable,colback=teal!4,colframe=teal!35!black,title={Before you move on}]
What does \gu{ઉત્તર} mean? (\textquotedblleft{}The north\textquotedblright{}.) Say it once more.
\end{tcolorbox}
\section[\texorpdfstring{\gu{ઊલટું} (ūlṭũ)}{ઊલટું (ūlṭũ)}]{\gu{ઊલટું} (ūlṭũ) — opposite; the wrong way round}

\label{lesson:GU-C414-ultu}



\begin{quote}
Before the new one: say the Gujarati for the west, then the Gujarati for the north.
\end{quote}

\subsection*{You'll want to know: \gu{ઊલટું}}
\textbf{\gu{ઊલટું}} — \emph{ūlṭũ} — \textquotedblleft{}opposite; the wrong way round\textquotedblright{}.

\begin{quote}
\gu{તમે} \gu{ઊલટી} \gu{દિશામાં} \gu{જાઓ} \gu{છો}.
\end{quote}

\begin{quote}
\emph{tame ūlṭī diśāmā̃ jāo cho.}
\end{quote}

\begin{quote}
You are going in the opposite direction.
\end{quote}

\begin{cousinweb}
\textbf{\gu{ઊલટું}} (\emph{ūlṭũ}) belongs with the verb CDIAL lists for Gujarati, \textquotedblleft{}to be the reverse of\textquotedblright{}, from \emph{*ullaṭati}, \textquotedblleft{}turns over\textquotedblright{}, a form it reconstructs (CDIAL 2368).
\end{cousinweb}

\begin{grammarlens}[title={What you've built}]
One more word for getting around.
\end{grammarlens}

\subsection*{Your turn}
\begin{itemize}
\raggedright
  \item \emph{Say it:} \emph{ūlṭũ}
  \item \emph{Say it:} \emph{ūlṭũ}, once more
  \item \emph{Say it:} \emph{tame ūlṭī diśāmā̃ jāo cho.} — You are going in the opposite direction.
  \item \emph{Recall:} say the Gujarati for the west, then the Gujarati for the north, then say \emph{ūlṭũ} again
\end{itemize}

\begin{tcolorbox}[breakable,colback=teal!4,colframe=teal!35!black,title={Before you move on}]
What does \gu{ઊલટું} mean? (\textquotedblleft{}Opposite; the wrong way round\textquotedblright{}.) Say it once more.
\end{tcolorbox}
\section[\texorpdfstring{\gu{હોકાયંત્ર} (hokāyantra)}{હોકાયંત્ર (hokāyantra)}]{\gu{હોકાયંત્ર} (hokāyantra) — a compass}

\label{lesson:GU-C414-hokayantra}



\begin{quote}
Before the new one: say the Gujarati for the north, then the Gujarati for the wrong way round.
\end{quote}

\subsection*{You'll want to know: \gu{હોકાયંત્ર}}
\textbf{\gu{હોકાયંત્ર}} — \emph{hokāyantra} — \textquotedblleft{}a compass\textquotedblright{}.

\begin{quote}
\gu{હોકાયંત્રની} \gu{સોય} \gu{ઉત્તર} \gu{બતાવે} \gu{છે}.
\end{quote}

\begin{quote}
\emph{hokāyantranī soy uttar batāve che.}
\end{quote}

\begin{quote}
The compass needle shows north.
\end{quote}

A neuter noun.

\begin{cousinweb}
\textbf{\gu{હોકાયંત્ર}} (\emph{hokāyantra}) is built on Sanskrit \emph{yantra}, \textquotedblleft{}an instrument, a device\textquotedblright{} (Monier-Williams); where its first part comes from is not settled.
\end{cousinweb}

\begin{grammarlens}[title={What you've built}]
One more word for getting around.
\end{grammarlens}

\subsection*{Your turn}
\begin{itemize}
\raggedright
  \item \emph{Say it:} \emph{hokāyantra}
  \item \emph{Say it:} \emph{hokāyantra}, once more
  \item \emph{Say it:} \emph{hokāyantranī soy uttar batāve che.} — The compass needle shows north.
  \item \emph{Recall:} say the Gujarati for the north, then the Gujarati for the wrong way round, then say \emph{hokāyantra} again
\end{itemize}

\begin{tcolorbox}[breakable,colback=teal!4,colframe=teal!35!black,title={Before you move on}]
What does \gu{હોકાયંત્ર} mean? (\textquotedblleft{}A compass\textquotedblright{}.) Say it once more.
\end{tcolorbox}
